\documentclass[12pt]{article}
\def\activityid{F07-X-v3}\usepackage[letterpaper,margin=.65in,top=.60in,bottom=.65in]{geometry}
\usepackage[T1]{fontenc}\usepackage[default]{sourcesanspro}
\usepackage{microtype,tikz,array,tabularx,fancyhdr,amsmath}\usepackage[hidelinks]{hyperref}
\usetikzlibrary{arrows.meta,calc}
\setlength{\parindent}{0pt}\setlength{\parskip}{7pt}\setlength{\headheight}{0pt}\setlength{\footskip}{26pt}\setlength{\emergencystretch}{2em}
\pagestyle{fancy}\fancyhf{}\renewcommand{\headrulewidth}{0pt}\renewcommand{\footrulewidth}{.4pt}
\fancyfoot[L]{\scriptsize Bellingham Math Circle / Week 7 / \activityid}\fancyfoot[R]{\scriptsize \thepage}
\definecolor{ink}{gray}{.12}\definecolor{soft}{gray}{.95}\color{ink}
\newcommand{\PageTitle}[2]{{\fontsize{10}{12}\selectfont\bfseries WEEK 7\quad GAME LAB\hfill #1\par}\vspace{3pt}{\fontsize{26}{29}\selectfont\bfseries #2\par}\vspace{2pt}}
\newcommand{\Subhead}[1]{\par\vspace{3pt}{\large\bfseries #1}\par}
\newcommand{\RuleBox}[1]{\par\begingroup\setlength{\fboxsep}{8pt}\colorbox{soft}{\parbox{\dimexpr\linewidth-16pt\relax}{#1}}\endgroup\par}
\newcommand{\WriteBox}[2]{\par\textbf{#1}\par\vspace{2pt}\begin{tikzpicture}\draw[black!40,line width=.5pt] (0,0) rectangle (\dimexpr\linewidth-.5pt\relax,#2);\end{tikzpicture}\par}
\newcommand{\Code}[1]{\texttt{#1}}

\newcommand{\StudentHeader}[1]{{\fontsize{10}{12}\selectfont\bfseries Week 7 / Strategy lab / #1\par}\vspace{10pt}}
\newcommand{\Problem}[2]{\par\vspace{4pt}\textbf{Problem #1:} #2\par}
\newcommand{\RecordBox}[2]{\par #1\par\vspace{2pt}\begin{tikzpicture}\draw[black!40,line width=.5pt](0,0)rectangle(\dimexpr\linewidth-.5pt\relax,#2);\end{tikzpicture}\par}

\hypersetup{pdftitle={Week 7: Extra / Grades 6--7}}
\begin{document}
\StudentHeader{Extra / Grades 6--7}
\Problem{1}{Use two piles. On a turn, take 1, 3, or 4 counters from exactly one pile; do not take more than that pile holds. Players alternate. Taking the very last counter wins. Play each starting pair below. Record the two pile sizes after every move, writing (left pile, right pile).}
\begin{center}\renewcommand{\arraystretch}{3.0}\begin{tabular}{c|p{3.1in}|p{1.6in}}Start & Pairs after each move & Winner: first or second\\\hline
(1,1) &  & \\
(1,3) &  & \\
(1,4) &  & \\
(4,4) &  & \\\end{tabular}\end{center}
\Problem{2}{Play again from (1,4). Try every legal first move from either pile. Reset before each attempt. Record which moves leave a position your opponent cannot win against your replies.}
\RecordBox{First moves and the pairs they leave:}{1.4in}

\newpage
\StudentHeader{Extra / Grades 6--7}
\Problem{3}{Give each single pile a number label. Set g(0)=0. For a larger pile, list the labels of every pile reachable by taking 1, 3, or 4. Its g-label is the smallest nonnegative integer missing from that list. For pile 1, the reachable label is 0, so g(1)=1. For pile 2, it is 1, so g(2)=0. Build these small cases, then complete the chart in order.}
\begin{center}\renewcommand{\arraystretch}{1.75}\begin{tabular}{c|p{1.8in}|p{1.8in}|c}Pile & Legal next piles & Their g-labels & Smallest missing\\\hline
0 & none & none & 0\\
1 & 0 & 0 & 1\\
2 & 1 & 1 & 0\\
3 &  &  & \\
4 &  &  & \\
5 &  &  & \\
6 &  &  & \\
7 &  &  & \\
8 &  &  & \\
9 &  &  & \\
10 &  &  & \\\end{tabular}\end{center}
\Problem{4}{Choose piles 4 and 5. Build every legal move from each. Use your chart to record the labels reached. Circle any nonnegative label smaller than the starting label that you can reach.}
\RecordBox{Starting labels and reachable labels:}{0.9in}

\newpage
\StudentHeader{Extra / Grades 6--7}
\Problem{5}{Use the number labels to revisit two-pile games. For each pair below, write both labels. If they differ, try to make them equal in one legal move. Record a move that does this, then test it in a game.}
\begin{center}\renewcommand{\arraystretch}{2.3}\begin{tabular}{c|p{1.1in}|p{2in}|p{1.2in}}Piles & Two g-labels & Move, if possible & New pair\\\hline
(1,3) &  &  & \\
(4,6) &  &  & \\
(5,6) &  &  & \\
(3,4) &  &  & \\\end{tabular}\end{center}
\Problem{6}{Start from (4,6), whose two labels are equal. List every legal first move and the labels of the new pair. For each new pair, find a reply that restores equal labels.}
\begin{center}\renewcommand{\arraystretch}{1.65}\begin{tabular}{p{1.5in}|p{2.4in}|p{1.8in}}First move & New pair and labels & My reply\\\hline
 &  & \\
 &  & \\
 &  & \\
 &  & \\
 &  & \\
 &  & \\\end{tabular}\end{center}

\end{document}
